% 全部内容为虚构排版示例，请按本人的实际答辩材料替换。
% 匿名评阅人的姓名、职称、单位和学科不公开时，可以不填写评阅人。
\begin{resolutiontext}
\begin{hitcommittee}
  % 非匿名评阅才使用下面的命令；不要填写猜测的姓名。
  % \hitreviewerrow{示例评阅人}{教授（是）}{示例大学}{示例学科}
  \hitcommitteerow{主席}{李教授}{教授（是）}{示例大学示例学院}{示例学科}
  \hitcommitteerow{委员}{示例委员甲}{教授（是）}{示例研究单位}{示例学科}
  \hitcommitteerow{委员}{示例委员乙}{教授（是）}{示例大学示例学院}{示例学科}
  \hitcommitteerow{秘书}{示例秘书}{副教授（否）}{示例大学示例学院}{示例学科}
\end{hitcommittee}

\noindent\textbf{答辩委员会决议：}

本段为虚构示例，仅用于展示原生排版。王同学的示例学位论文讨论了一个示例研究问题；论文的研究结果、评价意见、投票结果和学位授予建议，应以实际通过并签署的答辩决议为准。

原生模式的决议正文可以正常续页。请完整填写实际决议，不要为了控制页数删去有效意见。
\end{resolutiontext}
